% Figura 1: esquema da validação cruzada 5x2 (Dietterich, 1998; Alpaydin, 1999).
% Somente estrutura; nenhum valor numérico.
\documentclass[border=3pt]{standalone}
% Estilo comum das figuras esquemáticas (TikZ).
% Paleta Okabe--Ito (segura para daltonismo). A cor nunca é o único código:
% todo bloco de treino/teste também traz o rótulo textual.
\usepackage[T1]{fontenc}
\usepackage[utf8]{inputenc}
\usepackage[brazilian]{babel}
\usepackage{lmodern}
\usepackage{amsmath,amssymb}
\usepackage{tikz}
\usetikzlibrary{positioning,calc,decorations.pathreplacing,arrows.meta,fit,backgrounds}

\definecolor{cTreino}{HTML}{0072B2}   % azul
\definecolor{cTeste}{HTML}{E69F00}    % laranja
\definecolor{cModA}{HTML}{009E73}     % verde-azulado
\definecolor{cModB}{HTML}{CC79A7}     % púrpura-avermelhado
\definecolor{cAlerta}{HTML}{D55E00}   % vermelhão
\definecolor{cTinta}{HTML}{222222}
\definecolor{cCinza}{HTML}{6B6B6B}

\tikzset{
  every picture/.style={font=\sffamily\small, line cap=round, line join=round, text=cTinta},
  treino/.style={draw=cTreino!80!black, fill=cTreino!22, line width=0.5pt},
  teste/.style={draw=cTeste!70!black, fill=cTeste!45, line width=0.5pt},
  caixa/.style={draw=cTinta, rounded corners=2pt, line width=0.6pt, align=flush center,
                inner sep=4pt, fill=white},
  caixaA/.style={caixa, draw=cModA!80!black, fill=cModA!14},
  caixaB/.style={caixa, draw=cModB!80!black, fill=cModB!18},
  caixaT/.style={caixa, draw=cTreino!80!black, fill=cTreino!12},
  caixaTe/.style={caixa, draw=cTeste!70!black, fill=cTeste!25},
  caixaS/.style={caixa, draw=cCinza, fill=black!4},
  seta/.style={-{Stealth[length=2.4mm]}, line width=0.6pt, draw=cTinta},
  nota/.style={font=\sffamily\footnotesize, text=cCinza, align=flush left},
}

\begin{document}
\begin{tikzpicture}[x=1cm,y=1cm]
  \def\bh{0.62}   % altura das barras
  \def\dy{1.12}   % espaçamento vertical entre repetições
  % Cabeçalho
  \node[anchor=west, font=\sffamily\small\bfseries] at (-0.1,2.05)
    {Conjunto $\mathcal{D}$ com $n$ exemplos: a cada repetição $\ell$, uma permutação estratificada $\pi_\ell$ o divide em duas metades};
  \node[anchor=west, nota] at (-0.1,1.62)
    {$S_\ell^{(1)}$ e $S_\ell^{(2)}$, disjuntas e com $n/2$ exemplos cada; as metades mudam de uma repetição para outra.};
  % Legenda de cores
  \draw[treino] (12.75,1.49) rectangle ++(0.42,0.26);
  \node[anchor=west, font=\sffamily\footnotesize] at (13.2,1.62) {treino};
  \draw[teste]  (14.2,1.49) rectangle ++(0.42,0.26);
  \node[anchor=west, font=\sffamily\footnotesize] at (14.65,1.62) {teste};
  % Títulos das colunas
  \node[font=\sffamily\small\bfseries] at (4.75,1.0) {Dobra $j=1$};
  \node[nota, align=flush center] at (4.75,0.6) {treina em $S_\ell^{(1)}$, testa em $S_\ell^{(2)}$};
  \node[font=\sffamily\small\bfseries] at (10.4,1.0) {Dobra $j=2$ (papéis invertidos)};
  \node[nota, align=flush center] at (10.4,0.6) {treina em $S_\ell^{(2)}$, testa em $S_\ell^{(1)}$};
  \node[font=\sffamily\small\bfseries, align=flush center] at (14.55,0.8) {Variância\\[-1pt]da repetição};

  \foreach \r in {1,...,5} {
    \pgfmathsetmacro{\y}{-(\r-1)*\dy}
    % rótulo da repetição
    \node[anchor=west, font=\sffamily\small\bfseries] at (0,\y) {$\ell=\r$};
    \node[anchor=west, nota] at (0,\y-0.36) {partição $\pi_{\r}$};
    % Dobra 1: S1 treino | S2 teste
    \draw[treino] (2.5,\y-\bh/2) rectangle (4.75,\y+\bh/2);
    \draw[teste]  (4.75,\y-\bh/2) rectangle (7.0,\y+\bh/2);
    \node at (3.625,\y) {Treino $S_{\r}^{(1)}$};
    \node at (5.875,\y) {Teste $S_{\r}^{(2)}$};
    \draw[seta] (7.05,\y) -- (7.45,\y);
    \node[anchor=west] (d\r1) at (7.45,\y) {$\Delta_{\r 1}$};
    % Dobra 2: S1 teste | S2 treino
    \draw[teste]  (8.15,\y-\bh/2) rectangle (10.4,\y+\bh/2);
    \draw[treino] (10.4,\y-\bh/2) rectangle (12.65,\y+\bh/2);
    \node at (9.275,\y) {Teste $S_{\r}^{(1)}$};
    \node at (11.525,\y) {Treino $S_{\r}^{(2)}$};
    \draw[seta] (12.7,\y) -- (13.1,\y);
    \node[anchor=west] (d\r2) at (13.1,\y) {$\Delta_{\r 2}$};
    % variância
    \node (s\r) at (14.55,\y) {$\hat\sigma_{\r}^{2}$};
  }
  % Destaque do numerador de Dietterich
  \draw[cAlerta, line width=1pt, rounded corners=2pt]
    ($(d11.south west)+(-0.02,-0.02)$) rectangle ($(d11.north east)+(0.02,0.02)$);

  % Chaves
  \draw[decorate, decoration={brace, amplitude=5pt, mirror}, line width=0.6pt]
    (2.5,-4*\dy-0.55) -- (13.75,-4*\dy-0.55)
    node[midway, below=6pt, align=flush center, font=\sffamily\small]
    {$5$ repetições $\times$ $2$ dobras $=10$ diferenças pareadas\\
     $\Delta_{\ell j}=M^{A}_{\ell j}-M^{B}_{\ell j}$, medidas no mesmo conjunto de teste};
  \draw[decorate, decoration={brace, amplitude=5pt, mirror}, line width=0.6pt]
    (14.1,-4*\dy-0.55) -- (15.0,-4*\dy-0.55)
    node[midway, below=6pt, align=flush center, font=\sffamily\small] {$5$\\variâncias};

  % Caixas das estatísticas
  \node[caixa, draw=cAlerta, anchor=north west, text width=7.15cm, align=flush left] at (0,-6.35)
    {\textbf{Teste $\tilde t$ 5$\times$2cv (Dietterich)}\\[1pt]
     numerador: apenas $\Delta_{11}$ (em destaque)\\
     denominador: $\sqrt{\tfrac15\sum_{\ell} \hat\sigma_\ell^{2}}$\quad referência: $t_{5}$};
  \node[caixa, anchor=north west, text width=7.15cm, align=flush left] at (8.15,-6.35)
    {\textbf{Teste $F$ combinado 5$\times$2cv (Alpaydin)}\\[1pt]
     numerador: $\sum_{\ell}\sum_{j} \Delta_{\ell j}^{2}$ (as 10 diferenças)\\
     denominador: $2\sum_{\ell} \hat\sigma_\ell^{2}$\quad referência: $F_{10,\,5}$};
  \node[anchor=north west, nota] at (0,-7.95)
    {Por algoritmo: $10$ treinamentos, cada um com $n/2$ exemplos ($20$ para o par $A,B$). Os dois algoritmos usam exatamente as mesmas metades.};
\end{tikzpicture}
\end{document}
